\chapter{项目概述}

\section{课题背景}

在嵌入式视觉与智能控制应用中，摄像头与云台组成的目标跟踪系统具有较强的综合性：系统既需要完成图像采集与目标检测，又需要结合执行机构完成实时控制，还需要在有限算力平台上保证稳定性和可演示性。本项目以课程提供的 OrangePi、两自由度云台、摄像头与舵机驱动板为基础，构建一个面向现场演示的实时目标跟踪原型系统。

与追求复杂通用目标检测模型不同，当前阶段项目采用“稳闭环优先”的设计策略，将目标限定为红色高对比标记物。这样可以降低视觉算法复杂度，把主要精力集中在硬件打通、实时闭环、状态管理、日志记录和远程展示上，保证系统能够在 OrangePi 端稳定运行并形成可量化、可答辩说明的工程结果。

\section{项目目标}

本项目的阶段性目标包括：

\begin{enumerate}
  \item 完成 OrangePi、摄像头、PCA9685 舵机驱动板和双舵机云台的硬件连接；
  \item 在 OrangePi 上完成 Python 运行环境、OpenCV、I2C 与舵机控制依赖配置；
  \item 实现摄像头实时采集和红色高对比目标检测；
  \item 实现双舵机云台控制，使目标尽量保持在画面中心附近；
  \item 实现目标锁定、短时丢失、搜索重捕等状态管理逻辑；
  \item 实现运行日志记录和远程网页监控，方便调试、展示与后续数据分析；
  \item 初步完成 C/C++ 图像处理功能块，为后续 wheel 封装和性能对比实验准备。
\end{enumerate}

\section{小组成员}

本项目由三名成员合作完成，成员信息如表~\ref{tab:members} 所示。

\begin{table}[htbp]
  \centering
  \caption{小组成员信息}
  \label{tab:members}
  \begin{tabular}{lll}
    \toprule
    学号 & 姓名 & 当前主要工作 \\
    \midrule
    2352396 & {\namefont 禹尧珅} & 项目统筹、硬件调试、代码集成与报告整理 \\
    2352620 & 蒋文涛 & 视觉检测、跟踪流程与实验测试 \\
    2351283 & 吴凯 & 舵机控制、性能测试与展示材料整理 \\
    \bottomrule
  \end{tabular}
\end{table}

\section{当前完成情况}

截至本报告撰写时，系统已经完成了从硬件连接到远程网页监控的主流程打通。OrangePi 能够识别摄像头设备并读取图像帧，能够通过 I2C 总线识别 PCA9685 设备，并能够通过主程序控制两个舵机完成水平与俯仰运动。基础目标跟踪程序已可运行，网页监控模式也已跑通，可通过浏览器查看叠加后的实时画面。
